\documentclass[a4paper, 12pt]{article}
\usepackage[utf8]{inputenc}
\usepackage[T1]{fontenc}
\usepackage[french]{babel}
\usepackage{lmodern}
\usepackage{graphicx, float}
\usepackage{amsmath, amssymb, amsthm}
\usepackage{listings}
\usepackage[listings,skins]{tcolorbox}
\usepackage{xcolor}

\usepackage[hyphens]{url}
\usepackage[pdfauthor = {{Tristan Groult, Louis Guerain, Kibandou Bignang}}, pdftitle = {{Rapport projet Algorithmique des graphes}},pdfstartview = Fit, pdfpagelayout =
SinglePage, pdfnewwindow = true, bookmarksnumbered =
true, breaklinks, colorlinks, linkcolor = red, urlcolor =
black, citecolor = cyan, linktoc = all]{hyperref}

\usepackage[a4paper,margin=2.5cm]{geometry}

\renewcommand{\familydefault}{\sfdefault}
\definecolor{codegreen}{rgb}{0,0.6,0}
\definecolor{codegray}{rgb}{0.5,0.5,0.5}
\definecolor{codepurple}{rgb}{0.58,0,0.82}
\definecolor{codeblue}{rgb}{0.0, 0.4, 0.8}
\definecolor{darkWhite}{rgb}{0.90,0.90,0.90}

% Configuration pour les blocs de code (bash, python, etc.)
\lstdefinestyle{bigCode}{
    backgroundcolor=\color{darkWhite},
    commentstyle=\color{codegreen},
    keywordstyle=\color{codeblue},
    numberstyle=\tiny\color{codegray},
    stringstyle=\color{codepurple},
    basicstyle=\ttfamily\footnotesize,
    breakatwhitespace=false,
    breaklines=true,
    captionpos=b,
    keepspaces=true,
    showspaces=false,
    showstringspaces=false,
    showtabs=false,
    tabsize=2,
    language=bash,
    morekeywords={pip, install}
}

\lstset{style=bigCode}

\NewTotalTCBox{\commandbox}{ s v }
{verbatim,colupper=black,colback=darkWhite!75!white,colframe=white,left=0pt,boxsep=0px,right=0pt,top=2px,bottom=2px}
{\IfBooleanT{#1}{\textcolor{red}{\ttfamily\bfseries > }}%
\lstinline[language=bash,keywordstyle=\color{codeblue}\bfseries]^#2^}

\newcommand{\code}{\commandbox}

\title{Rapport projet d'Algorithmie des graphes \\ \large Résolution du problème du voyageur de commerce}
\author{Tristan GROULT, Louis GUERAIN, Kibandou BIGNANG}
\date{\today}

\begin{document}

\begin{figure}[t]
  \centering
  \begin{minipage}{0.3\textwidth}
    \centering
    \IfFileExists{logo_univ.png}{%
      \includegraphics[width=1\textwidth]{logo_univ.png}%
    }{%
      \fbox{\parbox[b][3cm][c]{\linewidth}{\centering logo\_univ.png manquant}}%
    }
  \end{minipage}
  \hfill
  \begin{minipage}{0.3\textwidth}
    \centering
    \IfFileExists{ufr_logo.png}{%
      \includegraphics[width=1\textwidth]{ufr_logo.png}%
    }{%
      \fbox{\parbox[b][3cm][c]{\linewidth}{\centering ufr\_logo.png manquant}}%
    }
  \end{minipage}
\end{figure}

\maketitle

\clearpage\setcounter{page}{2}

{
\hypersetup{hidelinks} % Sommaire en "noir"
\renewcommand{\contentsname}{Sommaire}
\tableofcontents % Affichage du sommaire ici
}

\clearpage

\section{Notice d'utilisation}

Ce projet dépend des bibliothèques \texttt{numpy} et d'autres modules standards.
Il dépend aussi de \texttt{flask} pour l'application web.

\begin{lstlisting}[caption=Installation des dépendances]
pip install numpy # (tout)
pip install flask # (web app)
\end{lstlisting}

Ou plus simplement :

\begin{lstlisting}
pip install requirement.txt
\end{lstlisting}

\section{Algorithme}

\subsection{Point le Plus Proche (PPP)}

\begin{tcolorbox}[colback=darkWhite!75!white,colframe=black!30,left=5pt,right=5pt,top=5pt,bottom=5pt]
\textbf{Algorithme : Procédure PPP (Point le Plus Proche)} \\
\textbf{Description :} Construit un cycle hamiltonien en insérant itérativement le point le plus proche du cycle actuel. \\
\textbf{Entrées :}
\begin{itemize}
    \item $L$ : Liste de points $P_1, \dots, P_n$
    \item $Q$ : Un point de départ choisi dans $L$
    \item $D$ : La matrice des distances $n \times n$
\end{itemize}
\textbf{Sorties :}
\begin{itemize}
    \item $c$ : Un cycle hamiltonien
    \item $l$ : La longueur du cycle
\end{itemize}
\textbf{Début :}
\begin{enumerate}
    \item $c \leftarrow$ cycle trivial composé du seul point $Q$
    \item $l \leftarrow 0$
    \item Retirer $Q$ de $L$
    \item Tant que $L$ n'est pas vide faire :
    \item \quad $(Q_i, Q_j) \leftarrow \texttt{TrouverCoupleMin}(L, c, D)$
    \item \quad $(c, l) \leftarrow \texttt{InsérerDansCycle}(c, Q_i, Q_j, D, l)$
    \item \quad Retirer $Q_i$ de $L$
    \item Retourner $c, l$
\end{enumerate}
\end{tcolorbox}

\vspace{5mm}

\begin{tcolorbox}[colback=darkWhite!75!white,colframe=black!30,left=5pt,right=5pt,top=5pt,bottom=5pt]
\textbf{Algorithme : InsérerDansCycle} \\
\textbf{Description :} Insère $Q_i$ dans $c$ soit entre $Q_j$ et son successeur, soit entre son prédécesseur et $Q_j$, selon le coût minimal, et met à jour la longueur du cycle. \\
\textbf{Entrées :}
\begin{itemize}
    \item $c$ : Cycle partiel
    \item $Q_i$ : Point à insérer
    \item $Q_j$ : Point d'ancrage dans $c$
    \item $D$ : Matrice des distances
    \item $l$ : Longueur actuelle du cycle
\end{itemize}
\textbf{Sorties :} $c$ (cycle mis à jour), $l$ (nouvelle longueur) \\
\textbf{Début :}
\begin{enumerate}
    \item Soit $Q_{next}$ le successeur de $Q_j$ dans $c$
    \item Soit $Q_{prev}$ le prédécesseur de $Q_j$ dans $c$
    \item $\text{coût}_{\text{droite}} \leftarrow D[Q_j, Q_i] + D[Q_i, Q_{next}] - D[Q_j, Q_{next}]$
    \item $\text{coût}_{\text{gauche}} \leftarrow D[Q_{prev}, Q_i] + D[Q_i, Q_j] - D[Q_{prev}, Q_j]$
    \item Si $\text{coût}_{\text{droite}} \leq \text{coût}_{\text{gauche}}$ alors :
    \item \quad Insérer $Q_i$ entre $Q_j$ et $Q_{next}$ dans $c$
    \item \quad $l \leftarrow l + \text{coût}_{\text{droite}}$
    \item Sinon :
    \item \quad Insérer $Q_i$ entre $Q_{prev}$ et $Q_j$ dans $c$
    \item \quad $l \leftarrow l + \text{coût}_{\text{gauche}}$
    \item Retourner $c, l$
\end{enumerate}
\end{tcolorbox}

\vspace{5mm}

\begin{tcolorbox}[colback=darkWhite!75!white,colframe=black!30,left=5pt,right=5pt,top=5pt,bottom=5pt]
\textbf{Algorithme : TrouverCoupleMin} \\
\textbf{Description :} Trouve le couple $(Q_i, Q_j)$ tel que $Q_i \in L$, $Q_j \in c$ et la distance $D[Q_i, Q_j]$ est minimale. \\
\textbf{Entrées :}
\begin{itemize}
    \item $L$ : Liste des points non encore visités
    \item $c$ : Le cycle partiel actuel
    \item $D$ : La matrice des distances $n \times n$
\end{itemize}
\textbf{Sorties :} $Q_i$ (point sélectionné dans $L$), $Q_j$ (point d'ancrage dans $c$) \\
\textbf{Début :}
\begin{enumerate}
    \item $d_{min} \leftarrow \infty$
    \item $Q_i \leftarrow \text{vide}$
    \item $Q_j \leftarrow \text{vide}$
    \item Pour chaque point $u$ appartenant à $L$ faire :
    \item \quad Pour chaque point $v$ appartenant à $c$ faire :
    \item \quad \quad Si $D[u, v] < d_{min}$ alors :
    \item \quad \quad \quad $d_{min} \leftarrow D[u, v]$
    \item \quad \quad \quad $Q_i \leftarrow u$
    \item \quad \quad \quad $Q_j \leftarrow v$
    \item Retourner $Q_i, Q_j$
\end{enumerate}
\end{tcolorbox}

\subsection{Point le Plus Proche Optimisé (OptPPP)}

\begin{tcolorbox}[colback=darkWhite!75!white,colframe=black!30,left=5pt,right=5pt,top=5pt,bottom=5pt]
\textbf{Algorithme : Procédure OptPPP} \\
\textbf{Description :} Améliore le cycle fourni par PPP en décroisant les arêtes (2-opt) tant que cela réduit la distance totale. \\
\textbf{Entrées :}
\begin{itemize}
    \item $c$ : Cycle hamiltonien initial (liste de points $P_1, \dots, P_n$)
    \item $D$ : Matrice des distances
\end{itemize}
\textbf{Sorties :} $c$ (cycle optimisé) \\
\textbf{Début :}
\begin{enumerate}
    \item $n \leftarrow \text{taille}(c)$
    \item $\text{amélioration} \leftarrow \text{Vrai}$
    \item Tant que $\text{amélioration}$ faire : \textit{// Répéter tant qu'on trouve une amélioration}
    \item \quad $\text{amélioration} \leftarrow \text{Faux}$
    \item \quad Pour $i$ allant de $1$ à $n-2$ faire : \textit{// Récupération de la première arête en provyant de la place pour la seconde}
    \item \quad \quad Pour $j$ allant de $i+2$ à $n$ faire : \textit{// Récupération de la seconde arête apès la première mais sans être adjacent à la première}
    \item \quad \quad \quad $u \leftarrow c[i]$
    \item \quad \quad \quad $v \leftarrow c[i+1]$
    \item \quad \quad \quad $x \leftarrow c[j]$
    \item \quad \quad \quad Si $j = n$ alors : \textit{// Ne pas dépasser n le dernier sommet}
    \item \quad \quad \quad \quad $y \leftarrow c[1]$
    \item \quad \quad \quad Sinon :
    \item \quad \quad \quad \quad $y \leftarrow c[j+1]$
    \item \quad \quad \quad $d_{actuelle} \leftarrow D[u, v] + D[x, y]$
    \item \quad \quad \quad $d_{nouvelle} \leftarrow D[u, x] + D[v, y]$
    \item \quad \quad \quad Si $d_{nouvelle} < d_{actuelle}$ alors : \textit{// le décroisement est intéressant}
    \item \quad \quad \quad \quad $p \leftarrow i+1$
    \item \quad \quad \quad \quad $q \leftarrow j$
    \item \quad \quad \quad \quad Tant que $p < q$ faire : \textit{// On inversse le sens des sommet dans le cycle}
    \item \quad \quad \quad \quad \quad Échanger $c[p]$ et $c[q]$
    \item \quad \quad \quad \quad \quad $p \leftarrow p + 1$
    \item \quad \quad \quad \quad \quad $q \leftarrow q - 1$
    \item \quad \quad \quad \quad $\text{amélioration} \leftarrow \text{Vrai}$
    \item Retourner $c$
\end{enumerate}
\end{tcolorbox}

\subsection{Arbre couvrant de poids minimum (OptPrim)}

\textit{Section à compléter.}

\subsection{Évaluation et séparation progressive ("branch and bound") (HDS)}

\textit{Section à compléter.}

\end{document}